\documentclass[11pt]{article}

% write-latex compliant: no fontspec, no XeLaTeX-only packages.
\usepackage[utf8]{inputenc}
\usepackage[T1]{fontenc}
\usepackage{lmodern}
\usepackage[margin=1in]{geometry}
\usepackage{amsmath, amssymb, amsthm, mathtools}
\usepackage{longtable}
\usepackage{booktabs}
\usepackage{enumitem}
\usepackage{hyperref}
\hypersetup{colorlinks=true, linkcolor=blue, urlcolor=blue, citecolor=blue}

\title{Tour --- Communication-Efficient Learning of Deep Networks from Decentralized Data}
\author{ McMahan et al., 2017 \\ \small Study repo: \href{ https://github.com/pleyva2004/communication-efficient-learning-of-deep-networks }{ pleyva2004/communication-efficient-learning-of-deep-networks } }
\date{}

\begin{document}
\maketitle

\noindent
A self-contained walking tour of this study. Read top-to-bottom for a complete first pass; jump by section if you already know the foundations. Estimated total time: \textbf{ 90 minutes}.

\iffalse
NOTE: this file is generated by the study-paper skill, Stage 7.
Mustache-style placeholders {{...}} are replaced per-paper before compilation.
Comments belong inside this iffalse block to keep them out of the rendered PDF.
\fi

\section*{1. Reader's contract}

By the end of this tour you will be able to:
\begin{itemize}[leftmargin=*]
  \item Derive why FedAvg cuts communication rounds and where non-IID data breaks it
  \item Reproduce the headline tradeoff and the Figure-1 barrier from runnable witnesses
  \item Map each proposed extension to the exact gap it closes (drift, bias, plateau, shared-init)
\end{itemize}

\noindent
This tour assumes comfort with multivariable calculus / gradients (refresh: \href{ https://github.com/pleyva2004/math-foundations/blob/main/concepts/gradient.md }{link}) and basic probability / expectation (refresh: \href{ https://github.com/pleyva2004/math-foundations/blob/main/concepts/expectation.md }{link}).

\noindent
Out of scope: differential privacy / secure aggregation; convergence theory for $E>1$.

\section*{2. Foundations walk}

Each row is one foundation node. Links resolve to the canonical concept page in the shared \href{https://github.com/pleyva2004/math-foundations}{math-foundations} repo.

\begin{longtable}{p{0.22\textwidth} p{0.45\textwidth} p{0.25\textwidth}}
\toprule
\textbf{Concept} & \textbf{Why it matters here} & \textbf{Link} \\
\midrule
\endhead
\bottomrule
\endfoot
expectation & defines the IID identity E[F_k]=f and all unbiasedness arguments & \href{https://github.com/pleyva2004/math-foundations/blob/main/concepts/expectation.md}{link} \\
empirical risk & the finite-sum objective f is the empirical risk being minimized & \href{https://github.com/pleyva2004/math-foundations/blob/main/concepts/empirical-risk.md}{link} \\
partition of a set & the client decomposition f=sum_k (n_k/n)F_k is exact because {P_k} partitions the data & \href{https://github.com/pleyva2004/math-foundations/blob/main/concepts/partition-of-a-set.md}{link} \\
unbiased estimator & IID means each F_k is an unbiased proxy for f; Horvitz-Thompson restores it under sampling & \href{https://github.com/pleyva2004/math-foundations/blob/main/concepts/unbiased-estimator.md}{link} \\
gradient & FedSGD aggregates per-client gradients into grad f & \href{https://github.com/pleyva2004/math-foundations/blob/main/concepts/gradient.md}{link} \\
gradient descent & FedSGD with C=1 is exactly full-batch GD & \href{https://github.com/pleyva2004/math-foundations/blob/main/concepts/gradient-descent.md}{link} \\
taylor theorem & the E>1 heterogeneity gap is a second-order Taylor remainder & \href{https://github.com/pleyva2004/math-foundations/blob/main/concepts/taylor-theorem.md}{link} \\
hessian & the gap is a covariance of local Hessians H_k with gradients & \href{https://github.com/pleyva2004/math-foundations/blob/main/concepts/hessian.md}{link} \\
covariance & the drift term is exactly Cov_k(H_k,g_k) & \href{https://github.com/pleyva2004/math-foundations/blob/main/concepts/covariance.md}{link} \\
convex function & Jensen safety of averaging needs convex F_k & \href{https://github.com/pleyva2004/math-foundations/blob/main/concepts/convex-function.md}{link} \\
jensens inequality & gives f(avg) <= avg f for convex objectives & \href{https://github.com/pleyva2004/math-foundations/blob/main/concepts/jensens-inequality.md}{link} \\
permutation group & hidden-unit relabelings are the symmetry behind the Fig-1 barrier & \href{https://github.com/pleyva2004/math-foundations/blob/main/concepts/permutation-group.md}{link} \\
control variates & the SCAFFOLD fix subtracts a control variate to cancel drift & \href{https://github.com/pleyva2004/math-foundations/blob/main/concepts/control-variates.md}{link} \\
horvitz thompson estimator & the unbiased aggregation fix is inverse-probability weighting & \href{https://github.com/pleyva2004/math-foundations/blob/main/concepts/horvitz-thompson-estimator.md}{link} \\
linear assignment problem & permutation alignment solves an assignment problem & \href{https://github.com/pleyva2004/math-foundations/blob/main/concepts/linear-assignment-problem.md}{link} \\
low rank factorization & federated LoRA communicates a low-rank adapter only & \href{https://github.com/pleyva2004/math-foundations/blob/main/concepts/low-rank-factorization.md}{link} \\
\end{longtable}

\section*{3. Paper concepts walk}

The paper graph. Read in topological order; each row links to a Markdown concept page (prose plus formal definition) and a Python witness (CPU-runnable, under 30 seconds).

\begin{longtable}{p{0.22\textwidth} p{0.45\textwidth} p{0.25\textwidth}}
\toprule
\textbf{Concept} & \textbf{Role in the paper} & \textbf{Files} \\
\midrule
\endhead
\bottomrule
\endfoot
01. Finite-Sum Objective & the objective everything minimizes & \texttt{01-finite-sum-objective} \\
02. Client-Partition Decomposition & rewrites the objective per client (exact) & \texttt{02-client-partition-decomposition} \\
03. IID vs Non-IID & defines the regime FedAvg must survive & \texttt{03-iid-noniid-dichotomy} \\
04. FedSGD as Exact Gradient Descent & the FedSGD baseline = exact GD & \texttt{04-fedsgd-gradient-descent} \\
05. Gradient- vs Model-Averaging Equivalence & why model-averaging = gradient-averaging (1 step) & \texttt{05-gradient-model-averaging-equivalence} \\
06. FedAvg: Iterated Local Steps & the FedAvg algorithm itself & \texttt{06-fedavg-local-iteration} \\
07. Local-Update Count u=nE/(KB) & the (C,E,B) compute-for-communication knob & \texttt{07-local-update-count} \\
08. Heterogeneity Gap = eta^2 Cov_k(H_k,g_k) & the leading E>1 deviation (client drift) & \texttt{08-heterogeneity-gap-covariance} \\
09. Client-Sampling Unbiasedness (Horvitz-Thompson) & justifies partial participation C<1 & \texttt{09-client-sampling-unbiasedness} \\
10. Aggregation Erratum: normalize by m_t & the published aggregation erratum & \texttt{10-aggregation-erratum} \\
11. Parameter Averaging & Jensen & why averaging is safe when convex & \texttt{11-parameter-averaging-jensen} \\
12. Shared-Init & the Permutation Barrier & why shared init is load-bearing & \texttt{12-shared-init-permutation-barrier} \\
\end{longtable}

\noindent\textbf{How the pieces fit.} Start from the finite-sum objective (01) and its exact client decomposition (02); the IID/non-IID dichotomy (03) says how faithful each client's slice is. FedSGD (04) is exact GD, and the one-step gradient/model-averaging equivalence (05) is the hinge FedAvg (06) generalizes by iterating local steps — quantified by the update count (07). The two deep facts are the heterogeneity gap (08), which is exactly why non-IID hurts, and the sampling/erratum pair (09,10), which fix the aggregation. Finally (11,12) explain when averaging parameters is safe and why a shared per-round initialization is non-negotiable.

\section*{4. Improvements walk}

Each proposal in \texttt{05-improvements.tex} gets exactly one validation file:
\begin{itemize}[leftmargin=*]
  \item \textbf{PROOF mode} --- a standalone LaTeX proof under \texttt{proofs/} with theorem statement, complete proof, and discussion of where it would break.
  \item \textbf{MEASUREMENT mode} --- an extended Python prototype under \texttt{improvements/} exposing \texttt{measure() -> dict} that prints a quantitative comparison table; deterministic under a fixed seed.
\end{itemize}

\begin{itemize}[leftmargin=*]
  \item \textbf{Control-Variate Drift Correction} --- PROOF --- \texttt{proofs/heterogeneity-control-variate.tex}
  \item \textbf{Unbiased Aggregation Estimators} --- PROOF --- \texttt{proofs/unbiased-aggregation.tex}
  \item \textbf{Adaptive Local-Work Schedule} --- MEASUREMENT --- \texttt{improvements/adaptive-local-work.py -> measure()}
  \item \textbf{Compute-Accounting Pareto (design)} --- link-only --- \texttt{design (E.2); see 05-improvements.tex}
  \item \textbf{Permutation-Aligned Averaging (Git Re-Basin for FL)} --- PROOF --- \texttt{proofs/permutation-aligned-averaging.tex}
  \item \textbf{Federated LoRA (design)} --- link-only --- \texttt{design (T.2); see 05-improvements.tex}
\end{itemize}

\section*{5. What to do next}

You have finished the tour. Pick one:
\begin{enumerate}[leftmargin=*]
  \item \textbf{Reproduce the sandbox claim.} \texttt{cd sandbox \&\& pip install -r requirements.txt \&\& python experiment.py}.
  \item \textbf{Run a measurement.} Pick any improvement marked MEASUREMENT, run \texttt{python improvements/<slug>.py}, change one knob, re-run.
  \item \textbf{Verify a proof.} Pick any improvement marked PROOF, audit the proof in \texttt{proofs/<slug>.tex}, then try to break it via the Discussion section.
  \item \textbf{Form an opinion.} Fill in \texttt{03-opinions.md}.
  \item \textbf{Extend the foundations graph.} Draft the missing concept at \href{https://github.com/pleyva2004/math-foundations}{pleyva2004/math-foundations}.
\end{enumerate}

\end{document}
